\subsectionlink{インフレーションの有効場理論}

有効場理論（EFT）は、素粒子物理学、さらには物性物理学においても、指導原理の役割を果たしてきた。EFTには、関心のあるエネルギースケールで重要な物理をまとめ上げる力がある。高エネルギーの、ほとんど関与しない物理の効果は、高次元演算子の係数に組み込まれる。これは、あるエネルギースケールでの現象論を探究するための方法である。これから行うのは、インフレーションの有効場理論を構築することである。その際、インフレーションをGoldstoneボソンの理論、すなわち時間並進のGoldstoneボソンの理論と捉えることができる。

\textbf{Goldstoneボソンの復習：}
Goldstoneボソンは素粒子物理学の至る所に現れる（これによってNambuは、まさに受けるに値する2008年のNobel賞を受賞した！）。メキシカンハット型のポテンシャルによって$\phi\to\langle\phi\rangle$となり、大域的な$U(1)$対称性$\phi\to e^{i\alpha}\phi$が自発的に破れる、最も単純な理論を考えよう。このとき、対称性を$\pi\to\pi+\alpha$という形で非線形に実現するGoldstoneボソン$\pi$が存在する。
\OriginalEquation{151}
したがって、場$\pi$の作用は、シフト対称性を備えた質量のないスカラー場の作用となる。
\OriginalEquation{152}
高階微分演算子は、高エネルギースケール$m/\lambda^{1/2}$の冪によって抑制される。

Goldstoneボソンの有名な例は、カイラル・ラグランジアンに現れるパイ中間子である。これは$SU(2)$カイラル・フレーバー対称性を非線形に実現するGoldstoneボソンであり、紫外（UV）側の理論であるQCDでは、クォークと反クォークの束縛状態に対応する。パイ中間子が創発的なスカラー場であることに注意しよう。QCDには基本的なスカラー場は存在しない。

\noindent\begin{minipage}{\linewidth}
\OriginalFigure{16}
\begin{center}
図16: $U(1)$対称性の自発的な破れを引き起こす、複素スカラー場のメキシカンハット・ポテンシャル。
\end{center}
\end{minipage}\par\medskip

\Needspace{4\baselineskip}
\textbf{Goldstoneボソンの理論としてのインフレーション：}
インフレーションのEFTをどのように構築すればよいだろうか。そのためには、インフレーションを最も本質的な観点から考える必要がある。インフレーションについて実際にわかっているのは、それが宇宙が準de Sitter状態にあった加速膨張の時期だということである。しかし、インフレーションは終わらなければならないので、厳密なde Sitter状態ではありえない。これは時間並進が自発的に破れていることを意味する。そこで、時間を測り、インフレーションを終わらせる物理的な時計が存在すると考えることにする。

この時計が何であろうと、一般相対論（GR）の座標不変性を利用すれば、この物理的な時計がゼロに設定される座標系へ移ることができる。任意の座標系から適切な時間微分同相変換を行い、時計の揺らぎがゼロとなる空間切片を選べばよい。例として、インフラトンが基本的なスカラー場であり（そう仮定しているわけではなく、単に例を挙げている）、$\delta\phi(\vec{x},t)\ne0$となる座標系にいるとしよう。このとき、次式が成り立つように時間微分同相変換$t\to\tilde{t}=t+\delta t(\vec{x},t)$を行うことができる（ここでは線形次数で示すが、任意の非線形次数へ一般化できる）。
\OriginalEquation{153}
さて、この座標系にいるとしよう。EFTの規則に従うことにする。その規則によれば、利用可能な自由度を用いて作用を書かなければならない。ここでは、それは計量の揺らぎだけである。揺らぎについて展開し、この問題の対称性と両立するすべての演算子を書き下さなければならない。今の場合、さまざまな空間切片の内部で空間座標を任意に変えることができ、それぞれの空間切片で異なる変換をしてよい。このことは、残るゲージ対称性が時間に依存する空間微分同相変換であることを意味する。
\OriginalEquation{154}

さらに、引き続きEFTの手順に従って、摂動について関心のある次数まで展開する（例えば、2点関数には2次、3点関数には3次、4点関数には4次、という具合である）。そのうえで、揺らぎの各次数において微分について展開する。高階微分項は、この問題のエネルギースケール$E$と何らかの高エネルギースケール$\Lambda$との比によって抑制される。

\noindent\begin{minipage}{\linewidth}
\OriginalFigure{17}
\begin{center}
図17: インフレーションを駆動する時計場が存在するならば、この時計を一様とする、特別な時間切片の取り方が存在する。
\end{center}
\end{minipage}\par\medskip

\subsubsection{ユニタリーゲージにおける作用の構成}

このユニタリーゲージで最も一般的なラグランジアンは何だろうか。ここでは文献\cite{Cheung:2007st}にほぼ沿って議論する。計量$g_{\mu\nu}$の関数であり、（線形に実現される）時間に依存する空間微分同相変換$x^i\to x^i+\xi^i(t,\vec{x})$のもとで不変な演算子を書き下さなければならない。時間に依存する空間微分同相変換は、実際、破れていない。言葉でいえば、インフレーションとは、時空の微分同相変換が時間に依存する空間微分同相変換へと自発的に破れた理論だということである\footnote[9]{微分同相変換はゲージの冗長性であることを覚えておこう（ゲージ対称性と呼ばれることもあるが、この呼び方はやや誤っている）。この点でインフレーションの理論は、$SU(2)\times U(1)$のゲージの冗長性が$U(1)$のゲージの冗長性へと自発的に破れる、低エネルギーでの素粒子物理学の標準模型に類似している。}。すべての微分同相変換に対して不変な、Riemannテンソルを含む通常の項に加えて、この系では対称性が減少しているため、多数の追加の項が許されるようになる。これらは、重力子に吸収された追加の自由度を記述する。例えば、$g^{00}$が空間微分同相変換のもとでスカラーであることは容易にわかる。したがって、ユニタリーゲージのラグランジアンには$g^{00}$が自由に現れてよい。
\OriginalEquation{155}
微分を含まない項は、$g^{00}$の多項式だけである。時空に特別な切片の取り方が存在するので、その切片を記述する幾何学的な量を書き下すことも許される。例えば、等時刻面の外部曲率$K_{\mu\nu}$は空間微分同相変換のもとでテンソルであり、作用に用いることができる。$n^\mu$を等時刻切片に直交するベクトルとすると、
\OriginalEquation{156}
である。ここで、$\nabla$は共変微分、$h_{\mu\nu}$は空間切片上の誘導計量であり、
\OriginalEquation{157}
と表される。

作用のどの項にも、一般の時間の関数を掛けてよいことに注意しよう。最も一般的なラグランジアンは、次のように書ける（証明は文献\cite{Cheung:2007st}の付録AおよびBを参照）。
\OriginalEquation{158}
ここで、省略記号は、揺らぎについてより高次の項、あるいはより多くの微分を含む項を表す。また、$\delta g^{00}=g^{00}+1$である。$\delta K_{\mu\nu}$は、等時刻面の外部曲率の、摂動のないFRWに対する変化を表し、$\delta K_{\mu\nu}=K_{\mu\nu}-a^2Hh_{\mu\nu}$である。ここで$h_{\mu\nu}$は誘導された空間計量である。上の作用において、選んだFRW解のまわりの線形摂動を含むのは最初の3項だけであり、その他はすべて明示的に2次以上であることに注意しよう。したがって、係数$c(t)$と$\Lambda(t)$は、与えられたFRWの発展$H(t)$をもつという要請、すなわち、この解のまわりでタドポール項が相殺されるという要請によって決まる。これらの係数を決定する前に、この簡略化が自明でないことを理解しておくことが重要である。背景解のまわりの1次で寄与する演算子は無限に多数あると考えられるだろう。しかし、$\delta K_{\mu\nu}$や$g^{00}+1$のような線形項の多項式として作用を書けば、ある演算子が1次、2次、あるいはさらに高次のどこから始まるかが明らかになる。式\eqref{eq:158}にあるもの以外のすべての線形項は微分を含み、部分積分によって、ここで考えた3つの線形項と、高次の共変な項との組み合わせにできる。したがって、\emph{摂動のない発展の履歴によって$c(t)$と$\Lambda(t)$が決まり、異なる模型間の違いは高次の項に組み込まれる}と結論できる。

ここで、与えられたFRWの発展が解となることを要求して、線形項を決定できる。すでに述べたように、$c$と$\Lambda$に比例する項だけが、エネルギー運動量テンソル
\OriginalEquation{159}
に、摂動の0次で消えない寄与を与え、したがってEinstein方程式の右辺に寄与する。インフレーションの間に主として関心があるのは、平坦なFRW宇宙
\OriginalEquation{160}
なので、Friedmann方程式は次のようになる。
\OriginalEquation{161}
\OriginalEquation{162}

$c$と$\Lambda$について解くと、作用\eqref{eq:158}を次のように書き換えられる。
\OriginalEquation{163}
すでに述べたように、上の作用に含まれる演算子の係数は、すべて一般の時間依存性をもってよい。しかし、ここで関心があるのは、1 Hubble時間の間に$H$と$\dot H$が大きく変化しない解である。したがって、他のすべての演算子についても同じことが成り立つと仮定するのが自然である。この仮定のもとでは、ラグランジアンは近似的に時間並進不変となる\footnote[10]{時間シフトが厳密な対称性となる極限は、$\dot H\to0$となるため、注意して取らなければならない。これは、ゴースト・インフレーションの議論で見るように、Goldstoneの空間的な運動項が消えることを意味する。}。そのため、ループ効果によって生じる時間依存性は、小さな対称性の破れのパラメーターによって抑制される\footnote[11]{この対称性は、時間微分同相変換の破れとは無関係であることに注意しよう。この対称性が$\phi$を用いた記述でどのように現れるかを見るには、適切な場の再定義を行えば、常に$\dot\phi=\mathrm{const.}$と仮定できることに注目するとよい。このように選ぶと、シフト対称性$\phi\to\phi+\mathrm{const.}$から、ユニタリーゲージのラグランジアンの時間並進不変性が従う。この対称性と、$\phi$のラグランジアンの時間並進対称性は、背景によって対角部分群へと破れる。この残った対称性が、ユニタリーゲージのラグランジアンにおける時間シフトである。}。この仮定は特に都合がよい。というのも、係数の急速な時間依存性は、指数関数的な膨張によって生じる摩擦に打ち勝つことがあり、その場合、インフレーションが力学的なアトラクターではなくなりうるからである。力学的なアトラクターであることは、標準的なFRW宇宙論の一様性問題を解決するために必要である。

このアプローチが、スカラーモードだけでなく重力についても最も一般的なラグランジアンを記述していることを強調しておきたい。例えば、高エネルギーの効果は、Riemannテンソルの摂動$\delta R_{\mu\nu\rho\sigma}$を含む演算子に組み込まれる。これらの補正は微分について高次なので、以下では明示的には議論しない。

単純なインフレーション模型をこの言葉でどのように書くか、いくつか例を挙げよう。最小の運動項とスローロール・ポテンシャル$V(\phi)$をもつ模型は、ユニタリーゲージでは次のように書ける。
\OriginalEquation{164}
Friedmann方程式から$\dot\phi_0(t)^2=-2M_P^2\dot H$および$V(\phi(t))=M_{\mathrm{Pl}}^2(3H^2+\dot H)$が得られるので、この作用は、最初の3項以外をすべてゼロとした式\eqref{eq:163}の形であることがわかる。明らかに、これが厳密に成り立つことはありえない。他のすべての項はループ補正によって生成されるからである。それらの項には、この単純なスローロール・インフレーション模型に対する、高エネルギーの物理のありうるすべての効果が組み込まれている。

より一般的な場合には、それぞれの$\phi$に作用する微分が高々1つであるような、ありうるすべてのラグランジアンが含まれる。すなわち、$\mathcal{L}=P(X,\phi)$であり、ここで$X=g^{\mu\nu}\partial_\mu\phi\partial_\nu\phi$である。摂動のない解$\phi_0(t)$のまわりでは、
\OriginalEquation{165}
となる。これは明らかに上記の形であり、$M_n^4(t)=\dot\phi_0(t)^{2n}\partial^nP/\partial X^n$を$\phi_0(t)$において評価したものとなっている。外部曲率を含む項は、1つのスカラーに作用する微分を2つ以上含み、厳密なde Sitter極限$\dot H\to0$において決定的に重要となる。これらはゴースト・インフレーションと、この枠組みで発見される新しい模型を再現する。
